\documentclass[aspectratio=169,11pt]{beamer}

% Compile with pdfLaTeX, LuaLaTeX, or XeLaTeX.
\usepackage[T1]{fontenc}
\usepackage{lmodern}
\usepackage{amsmath,amssymb}
\usepackage{booktabs}
\usepackage{tabularx}
\usepackage{xcolor}
\usepackage{ragged2e}

% ---------- Visual system ----------
\definecolor{Ink}{HTML}{10243E}
\definecolor{Ocean}{HTML}{167D9A}
\definecolor{Sky}{HTML}{DDEFF3}
\definecolor{Orange}{HTML}{E98032}
\definecolor{Paper}{HTML}{F7F9FB}
\definecolor{Muted}{HTML}{607080}

\setbeamercolor{normal text}{fg=Ink,bg=white}
\setbeamercolor{structure}{fg=Ocean}
\setbeamercolor{frametitle}{fg=Ink,bg=white}
\setbeamercolor{title}{fg=Ink}
\setbeamercolor{subtitle}{fg=Muted}
\setbeamercolor{alerted text}{fg=Orange}
\setbeamercolor{block title}{fg=white,bg=Ocean}
\setbeamercolor{block body}{fg=Ink,bg=Sky!50}
\setbeamercolor{footline}{fg=Muted,bg=white}

\setbeamertemplate{navigation symbols}{}
\setbeamertemplate{itemize item}{\color{Ocean}$\blacktriangleright$}
\setbeamertemplate{itemize subitem}{\color{Ocean}$\bullet$}
\setbeamertemplate{enumerate item}{\color{Ocean}\insertenumlabel.}
\setbeamertemplate{blocks}[rounded][shadow=false]
\setbeamertemplate{frametitle}{%
  \vspace{0.35cm}\insertframetitle\par
  \vspace{0.10cm}{\color{Ocean}\rule{\textwidth}{0.8pt}}
}
\setbeamertemplate{footline}{%
  \leavevmode
  \hbox{%
    \begin{beamercolorbox}[wd=.82\paperwidth,ht=2.8ex,dp=1.2ex,leftskip=.55cm]{footline}%
      SF2956 Topological Data Analysis
    \end{beamercolorbox}%
    \begin{beamercolorbox}[wd=.18\paperwidth,ht=2.8ex,dp=1.2ex,rightskip=.55cm plus1fil]{footline}%
      \insertframenumber/\inserttotalframenumber
    \end{beamercolorbox}%
  }
}

\setbeamerfont{title}{size=\LARGE,series=\bfseries}
\setbeamerfont{subtitle}{size=\large}
\setbeamerfont{frametitle}{size=\Large,series=\bfseries}
\setbeamerfont{footline}{size=\scriptsize}
\setbeamerfont{block title}{series=\bfseries}

\newcommand{\term}[1]{\textcolor{Ocean}{\textbf{#1}}}
\newcommand{\caution}[1]{\textcolor{Orange}{\textbf{#1}}}
\newcommand{\source}[1]{\vfill{\scriptsize\color{Muted}Source: #1}}

% ---------- Metadata ----------
\title{Topological Analysis of Temporal Trust Structure}
\subtitle{A proposed study of the Bitcoin Alpha network}
\author{Xiaoye Ma, Shengkai Sun}
\institute{SF2956 Topological Data Analysis}
\date{Project summary, 2026}

\begin{document}

% 1
\begin{frame}[plain]
  \vspace{0.8cm}
  {\color{Ocean}\rule{2.1cm}{3pt}}\\[0.45cm]
  {\usebeamerfont{title}\color{Ink}\inserttitle\par}
  \vspace{0.22cm}
  {\usebeamerfont{subtitle}\color{Muted}\insertsubtitle\par}

  \vfill
  \begin{block}{Project idea}
    Compare temporal windows through Jaccard pseudometrics, persistent homology, and distances between Betti curves.
  \end{block}

  \vfill
  {\small\insertauthor\hfill\insertinstitute\\[0.12cm]
  \color{Muted}\insertdate}
  \note{开场控制在一分钟内。说明这是 proposed project，不是已经完成的结果。主问题是用户信任行为的多尺度结构是否随时间变化。}
\end{frame}

% 2
\begin{frame}{Motivation}
  \begin{columns}[T,onlytextwidth]
    \column{0.54\textwidth}
      \begin{itemize}
        \item Bitcoin Alpha users rated other users from $-10$ to $+10$.
        \item Each rating has a direction, a sign, a weight, and a timestamp.
        \item The timestamps allow us to study how trust behaviour changes over time.
      \end{itemize}

    \column{0.40\textwidth}
      \begin{block}{Why a topological approach?}
        Counts describe network size. TDA describes how groups and higher-order structures persist across distance scales.
      \end{block}
  \end{columns}

  \vspace{0.45cm}
  \centering
  {\large\color{Ink}
  ratings $\longrightarrow$ behavioural similarity $\longrightarrow$ multi-scale structure}
  \note{不要把项目说成 Bitcoin 价格预测，也不要声称可以识别诈骗者。强调 ordinary counts 和 multi-scale structure 回答不同问题。}
\end{frame}

% 3
\begin{frame}{Bitcoin Alpha dataset}
  \begin{columns}[T,onlytextwidth]
    \column{0.52\textwidth}
      \renewcommand{\arraystretch}{1.18}
      \begin{tabularx}{\linewidth}{@{}Xr@{}}
        \toprule
        \textbf{Property} & \textbf{Value} \\
        \midrule
        Users & 3,783 \\
        Ratings & 24,186 \\
        Rating range & $-10$ to $+10$ \\
        Positive ratings & about 93\% \\
        \bottomrule
      \end{tabularx}

    \column{0.42\textwidth}
      \term{One row of data}
      \[
        (\text{source},\text{target},\text{rating},\text{time})
      \]
      \vspace{0.25cm}
      \begin{itemize}
        \item directed
        \item weighted and signed
        \item temporal
      \end{itemize}
  \end{columns}

  \source{Stanford Network Analysis Project, Bitcoin Alpha trust weighted signed network.}
  \note{解释四个字段。A 给 B 的评分与 B 给 A 的评分不同。时间戳是项目能够比较不同阶段的关键。}
\end{frame}

% 4
\begin{frame}{Research question}
  \vspace{0.25cm}
  \begin{block}{Main question}
    How does the multi-scale topological structure of user trust behaviour change over time?
  \end{block}

  \vspace{0.25cm}
  \begin{itemize}
    \item Do users form groups with similar positive-trust behaviour?
    \item Do these groups merge differently in early and late windows?
    \item Do $H_0$ and $H_1$ signatures vary across time?
    \item Can temporal windows be grouped into distinct regimes?
  \end{itemize}

  \vspace{0.25cm}
  \caution{Exploratory claim:} clear groups or temporal regimes may not exist.
  \note{这一页主动说明我们不会预设 cluster 一定存在。没有明显分组也是有效结果。}
\end{frame}

% 5
\begin{frame}{Temporal windows and preprocessing}
  \begin{columns}[T,onlytextwidth]
    \column{0.62\textwidth}
      \begin{enumerate}
        \item Convert Unix timestamps to dates.
        \item Divide the data into six-month windows.
        \item Keep the latest rating for repeated source--target pairs.
        \item Define ratings above zero as positive trust.
        \item Retain users with at least three positive ratings.
        \item Keep at most 200 highly active users per window.
      \end{enumerate}

    \column{0.32\textwidth}
      \begin{block}{Course notation}
        \[
          I=\{X_d\mid d\in D\}
        \]
        $D$: temporal windows\\[0.15cm]
        $X_d$: active users in window $d$
      \end{block}
  \end{columns}

  \vspace{0.25cm}
  The same rules will be applied to every window.
  \note{解释半年窗口是初始选择。活跃筛选避免大量空集合，200 人上限控制 H1 计算量。参数可在初步检查后调整，但规则应统一。}
\end{frame}

% 6
\begin{frame}{Jaccard pseudometric}
  For user $u$ in window $d$, define the positive-trust neighbourhood
  \[
    T_d(u)=\{v\mid u\text{ positively rates }v\}.
  \]

  \begin{columns}[T,onlytextwidth]
    \column{0.55\textwidth}
      \[
        d_d(u,v)=1-\frac{|T_d(u)\cap T_d(v)|}{|T_d(u)\cup T_d(v)|}.
      \]
      A small distance means that two users positively rate similar people.

    \column{0.39\textwidth}
      \begin{block}{Small example}
        \[
          T(A)=\{D,E,F\}
        \]
        \[
          T(B)=\{D,E,G\}
        \]
        \[
          d(A,B)=1-\frac{2}{4}=0.5
        \]
      \end{block}
  \end{columns}

  \caution{Important:} a small distance does not mean that $A$ directly trusts $B$.
  \note{解释为什么是 pseudometric：两个不同用户可能信任完全相同的一组人，因此距离为零。方向信息先进入 T_d(u)，随后距离本身是对称的。}
\end{frame}

% 7
\begin{frame}{Connection to the course pipeline}
  \centering
  \small
  \begin{tabularx}{0.96\textwidth}{@{}>{\bfseries\color{Ocean}}p{0.29\textwidth}X@{}}
    \toprule
    Project step & Course concept \\
    \midrule
    Temporal windows & $I=\{X_d\mid d\in D\}$ \\
    Jaccard distance & Metric restriction and pseudometric \\
    Groups at scale $t$ & Partitions $X_d/(d_d)_t$ \\
    Multi-scale clustering & Single linkage and dendrogram \\
    Geometric construction & Vietoris--Rips filtration \\
    Topological summary & $H_0$, $H_1$, barcodes, and Betti curves \\
    Window comparison & $l_1$ distance and hierarchical clustering \\
    \bottomrule
  \end{tabularx}

  \vspace{0.35cm}
  \normalsize

  \note{不逐条证明。老师需要看到数据表示、距离、filtration、homology、piecewise-constant functions 都有清楚对应。}
\end{frame}

% 8
\begin{frame}{Multi-scale groups and $H_0$}
  \begin{columns}[T,onlytextwidth]
    \column{0.52\textwidth}
      \term{Single-linkage dendrogram}
      \begin{itemize}
        \item At small $t$, only very similar users merge.
        \item Increasing $t$ produces larger behavioural groups.
        \item Merge heights record dissimilarity between groups.
      \end{itemize}

    \column{0.42\textwidth}
      \term{Zero-dimensional homology}
      \begin{itemize}
        \item $H_0$ records connected components.
        \item A bar dies when one component merges into another.
        \item A slow decrease of $\beta_0(t)$ indicates dispersed behaviour.
      \end{itemize}
  \end{columns}

  \vspace{0.35cm}
  \begin{block}{Primary analysis}
    $H_0$ will be the main result because it directly describes how behavioural groups merge across scales.
  \end{block}
  \note{连接回用户关于“群体性”的理解。提醒：不是选一个 cut 后才看结果，而是研究整个尺度范围。}
\end{frame}

% 9
\begin{frame}{Vietoris--Rips filtration and $H_1$}
  At scale $t$, connect users when $d_d(u,v)\leq t$ and fill every clique with a simplex.
  \[
    \mathrm{VR}_{t_1}(X_d)\subseteq
    \mathrm{VR}_{t_2}(X_d)\subseteq\cdots,
    \qquad t_1\leq t_2.
  \]

  \begin{columns}[T,onlytextwidth]
    \column{0.47\textwidth}
      \term{Barcode interpretation}
      \begin{itemize}
        \item birth: a feature appears
        \item death: the feature disappears
        \item long bar: persistence over a wider scale range
      \end{itemize}

    \column{0.47\textwidth}
      \term{Role of $H_1$}
      \begin{itemize}
        \item detects loops in the similarity space
        \item serves as a secondary analysis
        \item weak or short-lived loops remain reportable
      \end{itemize}
  \end{columns}

  \caution{$H_1$ loops are not direct trust cycles in the original network.}
  \note{可用空三角形与 filled triangle 的区别解释 H1。不要给每个 H1 bar 强行配社会学含义。}
\end{frame}

% 10
\begin{frame}{Comparing temporal windows}
  Each window produces Betti curves
  \[
    \beta_k^d(t)=\dim H_k\!\left(\mathrm{VR}_t(X_d)\right),
    \qquad k\in\{0,1\}.
  \]

  Compare windows with an $l_1$ distance
  \[
    D_k(d_i,d_j)=\int_0^1
    \left|\widetilde{\beta}_k^{d_i}(t)-
    \widetilde{\beta}_k^{d_j}(t)\right|\,dt.
  \]

  \vspace{0.15cm}
  \centering
  {\large
  Betti curves $\longrightarrow$ distance matrix $\longrightarrow$ window dendrogram}

  \vspace{0.35cm}
  \begin{block}{Two different dendrograms}
    User dendrograms describe structure within a window. The final dendrogram compares the windows themselves.
  \end{block}
  \note{如果每个窗口人数不同，优先使用归一化 beta0。窗口之间不需要逐个匹配相同用户，而是比较统一尺度区间上的曲线。}
\end{frame}

% 11
\begin{frame}{Planned outputs and interpretation}
  \begin{columns}[T,onlytextwidth]
    \column{0.46\textwidth}
      \term{Outputs}
      \begin{itemize}
        \item activity over time
        \item representative user dendrograms
        \item $H_0$ and $H_1$ barcodes
        \item Betti curves across windows
        \item window-distance heatmap
        \item temporal-window dendrogram
      \end{itemize}

    \column{0.48\textwidth}
      \term{Possible interpretations}
      \begin{itemize}
        \item Faster $\beta_0$ decrease: users connect at smaller scales.
        \item Persistent $H_1$: more complex similarity structure.
        \item Separate window branches: possible temporal regimes.
        \item Similar curves: relative stability under this representation.
      \end{itemize}
  \end{columns}

  \vspace{0.35cm}
  These are hypotheses to examine, not current findings.
  \note{展示最终报告可能包含哪些图。口头使用 possible / we will test whether，避免 we found。}
\end{frame}

% 12
\begin{frame}{Scope, feasibility, and limitations}
  \begin{columns}[T,onlytextwidth]
    \column{0.46\textwidth}
      \term{Deliberate scope}
      \begin{itemize}
        \item one main distance: Jaccard
        \item one main linkage: single linkage
        \item homology in dimensions 0 and 1
        \item at most 200 users per window
        \item no higher-dimensional computation
      \end{itemize}

    \column{0.48\textwidth}
      \term{Main limitations}
      \begin{itemize}
        \item rating strength is reduced to positive or non-positive
        \item negative ratings do not enter the main distance
        \item window width and activity threshold affect the output
        \item Jaccard similarity does not imply direct mutual trust
        \item the analysis does not establish causality
      \end{itemize}
  \end{columns}
  \note{这一页用来说明项目可控。局限不是缺点清单，而是表明我们理解 representation 做了哪些取舍。}
\end{frame}

% 13
\begin{frame}{Proposed project and requested feedback}
  \begin{block}{Summary}
    The project represents each temporal window as a pseudometric space, computes persistent topological summaries, and clusters windows through distances between their Betti curves.
  \end{block}

  \vspace{0.25cm}
  \term{Questions for discussion}
  \begin{itemize}
    \item Is the Jaccard representation appropriate for this project?
    \item Are six-month windows a reasonable initial choice?
    \item Is a primary focus on $H_0$, with $H_1$ as secondary, sufficient?
    \item Should we compare raw or normalised Betti curves?
  \end{itemize}

  \vfill
  \centering
  {\Large\color{Ocean}\textbf{Thank you}}
  \note{结尾重述完整 pipeline，并主动向老师确认四个具体选择。不要问过于宽泛的“这个项目可以吗”。}
\end{frame}

% ---------- Backup slides ----------
\appendix

\begin{frame}[noframenumbering]{Backup: likely questions}
  \small
  \term{Why TDA instead of ordinary community detection?}\\
  TDA follows connected components and higher-order structure across all distance thresholds rather than returning only one partition.

  \vspace{0.35cm}
  \term{Why Jaccard distance?}\\
  Each user is represented by a finite set of positively rated users. Jaccard distance directly compares the overlap of these sets.

  \vspace{0.35cm}
  \term{Why is the user space a pseudometric space?}\\
  Two different users may have identical positive-trust neighbourhoods and therefore distance zero.

  \vspace{0.35cm}
  \term{What if no clear clusters appear?}\\
  We will report that the selected representation does not reveal strongly separated behavioural groups or temporal regimes.
\end{frame}

\begin{frame}[noframenumbering]{Backup: parameter checks}
  Before the main analysis, we will examine
  \begin{itemize}
    \item ratings and active users per six-month window
    \item the effect of minimum activity thresholds 2 and 3
    \item whether the 200-user cap removes a substantial share of activity
    \item sensitivity to nearby window widths if the default is too sparse
    \item raw versus normalised $\beta_0$ curves
  \end{itemize}

  \vspace{0.35cm}
  Parameter changes will be documented and applied consistently across windows.
\end{frame}

\begin{frame}[noframenumbering]{References}
  \small
  \begin{thebibliography}{9}
    \bibitem{snap}
      Stanford Network Analysis Project.
      \newblock \emph{Bitcoin Alpha trust weighted signed network}.
      \newblock \href{https://snap.stanford.edu/data/soc-sign-bitcoin-alpha.html}{snap.stanford.edu/data/soc-sign-bitcoin-alpha.html}

    \bibitem{kumar2016}
      S. Kumar, F. Spezzano, V. S. Subrahmanian, and C. Faloutsos.
      \newblock \emph{Edge Weight Prediction in Weighted Signed Networks}.
      \newblock IEEE International Conference on Data Mining, 2016.

    \bibitem{course}
      SF2956 Topological Data Analysis.
      \newblock Course notes and lecture handouts, 2026.
  \end{thebibliography}
\end{frame}

\end{document}
